% SEGMENT OLP-0406-S01
% SOURCE CORRECTION TA-MVL-027: correct the part metadata from first-order-logic to many-valued-logic.
% Part: many-valued-logic
% Chapter: sequent-calculus
% Section: propositional-rules

\documentclass[../../../include/open-logic-section]{subfiles}

\begin{document}

\olfileid{mvl}{seq}{prl}

\olsection{தேர்ந்தெடுக்கப்பட்ட தருக்கங்களுக்கான கூற்று விதிகள்}

ஓர்~$n$-பக்கத் தொடரணிக் கணியத்தில் ஓர் இணைப்பிக்கான அனுமான விதிகள், அந்த
இணைப்பியைப் பண்புறுத்தும் வாய்மைச் சார்பை மட்டுமே சார்ந்துள்ளன. ஆகவே சில
இணைப்பி வெவ்வேறு தருக்கங்களில் ஒரே வாய்மைச் சார்பால் வரையறுக்கப்பட்டால்,
அந்த இணைப்பிக்கான $n$-பக்கத் தொடரணி விதிகள் அத்தருக்கங்களில் ஒன்றே.

\subsection{$\lnot$ க்கான விதிகள்}

$\lnot$ க்கான பின்வரும் விதிகள் லூகாசியேவிச்~(\L ukasiewicz), க்ளீனி
தருக்கங்களுக்கும் அவற்றின் மாறுபாடுகளுக்கும் பொருந்துகின்றன.

\begin{defish}
  \begin{center}
\AxiomC{$ \Gamma \nSequent \Pi \nSequent \Delta, !A $}
\RightLabel{\iR{\lnot}{\False}}
\UnaryInfC{$ \lnot !A, \Gamma \nSequent \Pi \nSequent \Delta$}
\DisplayProof
\\[2ex]
\AxiomC{$ \Gamma \nSequent !A, \Pi \nSequent \Delta $}
\RightLabel{\iR{\lnot}{\Undef}}
\UnaryInfC{$\Gamma \nSequent \lnot !A, \Pi \nSequent \Delta$}
\DisplayProof
\\[2ex]
\AxiomC{$!A, \Gamma \nSequent \Pi \nSequent \Delta$}
\RightLabel{\iR{\lnot}{\True}}
\UnaryInfC{$\Gamma \nSequent \Pi \nSequent \Delta,  \lnot !A$}
\DisplayProof
  \end{center}
\end{defish}

$\lnot$ க்கான பின்வரும் விதிகள் கோடல்~(G\"odel) தருக்கத்திற்குப்
பொருந்துகின்றன.

\begin{defish}
\AxiomC{$ \Gamma \nSequent !A, \Pi \nSequent \Delta, !A $}
\RightLabel{\iR{\lnot}{\False}[\LogGod]}
\UnaryInfC{$ \lnot !A, \Gamma \nSequent \Pi \nSequent \Delta$}
\DisplayProof
\hfill
\AxiomC{$!A, \Gamma \nSequent \Pi \nSequent \Delta$}
\RightLabel{\iR{\lnot}{\True}[\LogGod]}
\UnaryInfC{$\Gamma \nSequent \Pi \nSequent \Delta,  \lnot !A$}
\DisplayProof
\hfill
\end{defish}

(கோடல்~(G\"odel) தருக்கத்தில், $\lnot !A$ ஆனது~$\Undef$ என்ற மதிப்பை
ஒருபோதும் பெற முடியாது; ஆகவே நடு இடத்திற்கான விதி இல்லை.)

\subsection{$\land$ க்கான விதிகள்}

இவை லூகாசியேவிச்~(\L ukasiewicz), வலு க்ளீனி, கோடல்~(G\"odel)
தருக்கங்களில்~$\land$ க்கான விதிகள்.

\begin{defish}
\begin{center}
  \AxiomC{$!A, !B, \Gamma \nSequent \Pi \nSequent \Delta$}
  \RightLabel{$\iR{\land}{\False}$}
  \UnaryInfC{$!A \land !B, \Gamma \nSequent \Pi \nSequent \Delta$}
  \DisplayProof
\\[2ex]
  \AxiomC{$\Gamma \nSequent !A, \Pi \nSequent !A, \Delta$}
  \AxiomC{$\Gamma \nSequent !B, \Pi \nSequent !B, \Delta$}
  \AxiomC{$\Gamma \nSequent !A, !B, \Pi \nSequent \Delta$}
  \RightLabel{$\iR\land\Undef$}
  \TrinaryInfC{$\Gamma \nSequent !A \land !B, \Pi \nSequent \Delta$}
  \DisplayProof
  \\[2ex]  
  \AxiomC{$\Gamma \nSequent \Pi \nSequent \Delta, !A$}
  \AxiomC{$\Gamma \nSequent \Pi \nSequent \Delta, !B$}
  \RightLabel{$\iR\land\True$}
  \BinaryInfC{$\Gamma \nSequent \Pi \nSequent \Delta, !A \land !B$}
  \DisplayProof
\end{center}
\end{defish}

\subsection{$\lor$ க்கான விதிகள்}

இவை லூகாசியேவிச்~(\L ukasiewicz), வலு க்ளீனி, கோடல்~(G\"odel)
தருக்கங்களில்~$\lor$ க்கான விதிகள்.

\begin{defish}
\begin{center}
 \AxiomC{$!A, \Gamma \nSequent \Pi \nSequent \Delta$}
  \AxiomC{$!B, \Gamma \nSequent \Pi \nSequent \Delta$}
  \RightLabel{$\iR\lor\False$}
  \BinaryInfC{$!A \lor !B, \Gamma \nSequent \Pi \nSequent \Delta$}
  \DisplayProof
\\[2ex]
  \AxiomC{$!A, \Gamma \nSequent !A, \Pi \nSequent \Delta$}
  \AxiomC{$!B, \Gamma \nSequent !B, \Pi \nSequent \Delta$}
  \AxiomC{$\Gamma \nSequent !A, !B, \Pi \nSequent \Delta$}
  \RightLabel{$\iR\lor\Undef$}
  \TrinaryInfC{$\Gamma \nSequent !A \lor !B, \Pi \nSequent \Delta$}
  \DisplayProof
  \\[2ex]  
  \AxiomC{$\Gamma \nSequent \Pi \nSequent \Delta, !A, !B$}
  \RightLabel{$\iR\lor\True$}
  \UnaryInfC{$\Gamma \nSequent \Pi \nSequent \Delta, !A \lor !B$}
  \DisplayProof
\end{center}
\end{defish}

\subsection{$\lif$ க்கான விதிகள்}

இவை லூகாசியேவிச்~(\L ukasiewicz) தருக்கத்தில்~$\lif$ க்கான விதிகள்.

\begin{defish}
\begin{center}
  \AxiomC{$\Gamma \nSequent \Pi \nSequent \Delta, !A$}
  \AxiomC{$!B, \Gamma \nSequent \Pi \nSequent \Delta$}
  \RightLabel{$\iR\lif\False[\LogLuk[3]]$}
  \BinaryInfC{$!A \lif !B, \Gamma \nSequent \Pi \nSequent \Delta$}
  \DisplayProof
  \\[2ex]
  \AxiomC{$\Gamma \nSequent !A, !B, \Pi \nSequent \Delta$}
  \AxiomC{$!B, \Gamma \nSequent \Pi \nSequent \Delta, !A$}
  \RightLabel{$\iR\lif\Undef[\LogLuk[3]]$}
  \BinaryInfC{$\Gamma \nSequent !A \lif !B, \Pi \nSequent \Delta$}
  \DisplayProof
\\[2ex]
  \AxiomC{$!A, \Gamma \nSequent !B, \Pi \nSequent \Delta, !B$}
  \AxiomC{$!A, \Gamma \nSequent !A, \Pi \nSequent \Delta, !B$}
  \RightLabel{$\iR{\lif}{\True}[\LogLuk[3]]$}
  \BinaryInfC{$\Gamma \nSequent \Pi \nSequent \Delta, !A \lif !B$}
  \DisplayProof
\end{center}
\end{defish}

இவை வலு க்ளீனி தருக்கத்தில்~$\lif$ க்கான விதிகள்.

\begin{defish}
\begin{center}
  \AxiomC{$\Gamma \nSequent \Pi \nSequent \Delta, !A$}
  \AxiomC{$!B, \Gamma \nSequent \Pi \nSequent \Delta$}
  \RightLabel{$\iR\lif\False[\LogKs]$}
  \BinaryInfC{$!A \lif !B, \Gamma \nSequent \Pi \nSequent \Delta$}
  \DisplayProof
  \\[2ex]
  \AxiomC{$!B, \Gamma \nSequent !B, \Pi \nSequent \Delta$}
  \AxiomC{$\Gamma \nSequent !A, !B, \Pi \nSequent \Delta$}
  \AxiomC{$\Gamma \nSequent !A, \Pi \nSequent \Delta, !A$}
  \RightLabel{$\iR\lif\Undef[\LogKs]$}
  \TrinaryInfC{$\Gamma \nSequent !A \lif !B, \Pi \nSequent \Delta$}
  \DisplayProof
\\[2ex]
  \AxiomC{$!A, \Gamma \nSequent \Pi \nSequent \Delta, !B$}
  \RightLabel{$\iR{\lif}{\True}[\LogKs]$}
  \UnaryInfC{$\Gamma \nSequent \Pi \nSequent \Delta, !A \lif !B$}
  \DisplayProof
\end{center}
\end{defish}

இவை கோடல்~(G\"odel) தருக்கத்தில்~$\lif$ க்கான விதிகள்.

\begin{defish}
\begin{center}
  \AxiomC{$\Gamma \nSequent !A, \Pi \nSequent \Delta, !A$}
  \AxiomC{$!B, \Gamma \nSequent \Pi \nSequent \Delta$}
  \RightLabel{$\iR\lif\False[\LogGod[3]]$}
  \BinaryInfC{$!A \lif !B, \Gamma \nSequent \Pi \nSequent \Delta$}
  \DisplayProof
  \\[2ex]
  \AxiomC{$\Gamma \nSequent !B, \Pi \nSequent \Delta$}
  \AxiomC{$\Gamma \nSequent \Pi \nSequent \Delta, !A$}
  \RightLabel{$\iR\lif\Undef[\LogGod[3]]$}
  \BinaryInfC{$\Gamma \nSequent !A \lif !B, \Pi \nSequent \Delta$}
  \DisplayProof
\\[2ex]
  \AxiomC{$!A, \Gamma \nSequent !B, \Pi \nSequent \Delta, !B$}
  \AxiomC{$!A, \Gamma \nSequent !A, \Pi \nSequent \Delta, !B$}
  \RightLabel{$\iR{\lif}{\True}[\LogGod[3]]$}
  \BinaryInfC{$\Gamma \nSequent \Pi \nSequent \Delta, !A \lif !B$}
  \DisplayProof
\end{center}
\end{defish}


\begin{sidewaysfigure}
  \begin{center}  
  \AxiomC{$A \nSequent A \nSequent A$}
  \RightLabel{\iR \Weakening \True}
  \UnaryInfC{$A \nSequent A \nSequent B, A$}
  \RightLabel{\iR \Weakening \Undef}
  \UnaryInfC{$A \nSequent B, A \nSequent B, A$}
  \RightLabel{\iR \Weakening \Undef}
  \UnaryInfC{$A \nSequent A, B, A \nSequent B, A$}
  \AxiomC{$A \nSequent A \nSequent A$}
  \RightLabel{\iR \Weakening \True}
  \UnaryInfC{$A \nSequent A \nSequent A, A$}
  \RightLabel{\iR \Weakening \True}
  \UnaryInfC{$A \nSequent A \nSequent B, A, A$}
  \RightLabel{\iR \Weakening \False}
  \UnaryInfC{$B, A \nSequent A \nSequent B, A, A$}
  \RightLabel{$\iR\lif\Undef$}
  \BinaryInfC{$A \nSequent A \lif B, A \nSequent B, A$}
  \AxiomC{$B \nSequent B \nSequent B$}
  \RightLabel{\iR \Weakening \Undef}
  \UnaryInfC{$B \nSequent A, B \nSequent B$}
  \RightLabel{\iR \Exchange \Undef}
  \UnaryInfC{$B \nSequent B, A \nSequent B$}
  \RightLabel{\iR \Weakening \Undef}
  \UnaryInfC{$B \nSequent A, B, A \nSequent B$}
  \RightLabel{\iR \Weakening \False}
  \UnaryInfC{$A, B \nSequent A, B, A \nSequent B$}
  \RightLabel{\iR \Exchange \False}
  \UnaryInfC{$B, A \nSequent A, B, A \nSequent B$}
  \AxiomC{$A \nSequent A \nSequent A$}
  \RightLabel{\iR \Weakening \True}
  \UnaryInfC{$A \nSequent A \nSequent B, A$}
  \RightLabel{\iR \Weakening \False}
  \UnaryInfC{$B, A \nSequent A \nSequent B, A$}
  \RightLabel{\iR \Weakening \False}
  \UnaryInfC{$B, B, A \nSequent A \nSequent B, A$}
  \RightLabel{$\iR\lif\Undef$}
  \BinaryInfC{$B, A \nSequent A \lif B, A \nSequent B$}
  \RightLabel{$\iR\lif\False$}
  \BinaryInfC{$A \lif B, A \nSequent A \lif B, A \nSequent B$}
  \DisplayProof
  \end{center}
  \caption{$\LogLuk[3]$ இலுள்ள !!{derivation} எடுத்துக்காட்டு}
\end{sidewaysfigure}

\end{document}
